\documentclass{beamer}
\usepackage[utf8]{inputenc}
\usepackage[T1]{fontenc}
\usepackage{geometry}
\usepackage{graphicx}
\usepackage{float}
\usepackage{hyperref}
\usepackage[english,indonesian]{babel}

\title{Optimasi Pemrosesan Template LaTeX Menggunakan Pendekatan Berbasis AI}
\author{Academic Support Agent}
\date{\today}

\begin{document}
\maketitle

\begin{abstract}

\end{abstract}

\chapter{PENDAHULUAN}


\chapter{TINJAUAN PUSTAKA}


% Original content (commented out):
% \documentclass[10pt,xcolor={dvipsnames}]{beamer}
\usetheme[
%%% option passed to the outer theme
%    progressstyle=fixedCircCnt,   % fixedCircCnt, movingCircCnt (moving is deault)
  ]{Feather}
  
% If you want to change the colors of the various elements in the theme, edit and uncomment the following lines

% Change the bar colors:
\setbeamercolor{Feather}{fg=NavyBlue!20,bg=NavyBlue}

% Change the color of the structural elements:
\setbeamercolor{structure}{fg=NavyBlue}

% Change the frame title text color:
\setbeamercolor{frametitle}{fg=black!5}

% Change the normal text colors:
\setbeamercolor{normal text}{fg=black!75,bg=gray!5}

%% Change the block title colors
\setbeamercolor{block title}{use=Feather,bg=Feather.fg, fg=black!90} 


% Change the logo in the upper right circle:
%\renewcommand{\logofile}{example-grid-100x100pt} 
%% This is an image that comes with the LaTeX installation
% Adjust scale of the logo w.r.t. the circle; default is 0.875
% \renewcommand{\logoscale}{0.55}

% Change the background image on the title and final page.
% It stretches to fill the entire frame!
% \renewcommand{\backgroundfile}{example-grid-100x100pt}

%-------------------------------------------------------
% INCLUDE PACKAGES
%-------------------------------------------------------

%\usepackage[utf8]{inputenc}
%\usepackage[english]{babel}
%\usepackage[T1]{fontenc}
% \usepackage{helvet}

%% Load different font packages to use different fonts
%% e.g. using Linux Libertine, Linux Biolinum and Inconsolata
% \usepackage{libertine}
% \usepackage{zi4}

%------------ สำหรับการใช้ฟ้อนท์ภาษาไทย -----------%
\usepackage[no-math]{fontspec}
\usepackage{xunicode}
\usepackage{xltxtra}
\usepackage{polyglossia}

%% e.g. using fonts Carlito and Caladea
\usepackage{carlito}
\usepackage{caladea}
\usepackage{zi4}

\setdefaultlanguage{english}
\setotherlanguage{thai} 
\defaultfontfeatures{Mapping=tex-text} % so that the quotation marks display correctly

\newfontfamily{\thaifont}[Scale=MatchLowercase,Mapping=tex-text,Script=Thai]{THSarabunNew_Bold.ttf}
%\newfontfamily{\thaifont}[Mapping=tex-text,Script=Thai]{Noto Serif Thai}
\usepackage[Latin,Thai]{ucharclasses}
\setTransitionTo{Thai}{\thaifont}
\setTransitionFrom{Thai}{\normalfont}
% % % % % % % % % % % % % % % % %
% คำสั่งสำหรับตัดคำ
\XeTeXlinebreaklocale "th_TH"
\XeTeXlinebreakskip=0pt plus 1pt
% % % % % % % % % % % % % %



%% e.g. using Venturis ADF Serif and Sans
% \usepackage{venturis}

%-------------------------------------------------------
% DEFFINING AND REDEFINING COMMANDS
%-------------------------------------------------------

% colored hyperlinks
\newcommand{\chref}[2]{
  \href{#1}{{\usebeamercolor[bg]{Feather}#2}}
}

%-------------------------------------------------------
% INFORMATION IN THE TITLE PAGE
%-------------------------------------------------------

\title[] % [] is optional - is placed on the bottom of the sidebar on every slide
{ % is placed on the title page
      \textbf{บีมเมอร์ธีมขนนก Feather}
}

\subtitle[บีมเมอร์ธีมขนนก Feather]
{
      \textbf{v. 1.1.0}
}

\author[Lilyana Vankova, LianTze Lim, B. Intiyot]
{      Lilyana Vankova \\
      {\ttfamily lilqna.v@gmail.com}\\[1em]
      with v1.1 modifications by LianTze Lim (Overleaf)\\
      with Thai translation by B.Intiyot
}

\institute[]
{%
      คณะคณิตศาสตร์, 	สนเทศศาสตร์และเทคโนโลยีสารสนเทศ\\
      มหาวิทยาลัย Plovdiv ``Paisii Hilendarski''
}

\date{\today}

%-------------------------------------------------------
% THE BODY OF THE PRESENTATION
%-------------------------------------------------------

\begin{document}

%-------------------------------------------------------
% THE TITLEPAGE
%-------------------------------------------------------

{\1% % this is the name of the PDF file for the background
\begin{frame}[plain,noframenumbering] % the plain option removes the header from the title page, noframenumbering removes the numbering of this frame only
  \titlepage % call the title page information from above
\end{frame}}


\begin{frame}{เนื้อหา}{}
\tableofcontents
\end{frame}

%-------------------------------------------------------
\section{บทนำ}
%-------------------------------------------------------
\subsection{ใบอนุญาต}
\begin{frame}{บทนำ}{ใบอนุญาต}
%-------------------------------------------------------

  \begin{itemize}
    \item<1-> รูปขนนก Feather ไม่อยู่ภายใต้กฏ copyright ฉันใช้รูปจาก \chref{http://www.vectors-free.com/}{http://www.vectors-free.com/} คุณสามารถใช้รูปขนนก Feather สำหรับวัตถุประสงค์อะไรก็ได้
    \item<2-> ส่วนอื่น ๆ ของธีมอยู่ภายใต้ใบอนุญาต GNU General Public License v. 3 (GPLv3) \chref{http://www.gnu.org/licenses/}{http://www.gnu.org/licenses/} นี่หมายความว่าคุณสามารถที่จะแจกจ่าย และ/หรือแก้ไขมันภายใต้ใบอนุญาตเดียวกัน 
  \end{itemize}
\end{frame}

%-------------------------------------------------------
\section{การติดตั้ง}
%-------------------------------------------------------
\subsection{ไฟล์ต้นกำเนิด}
\begin{frame}{การติดตั้ง}{ไฟล์ต้นกำเนิด}
%-------------------------------------------------------

\begin{block}{}
ธีมนี้มีไฟล์ต้นกำเนิด 4 ไฟล์คือ
  \begin{itemize}
    \item {\tt beamercolorthemeFeather.sty}
    \item {\tt beamerouterthemeFeather.sty}
    \item {\tt beamerinnerthemeFeather.sty}
    \item {\tt beamerthemeFeather.sty}
  \end{itemize}
\end{block}
\end{frame}

%-------------------------------------------------------
\subsection{การติดตั้งแบบเฉพาะที่และแบบครอบคลุม}
\begin{frame}[shrink]{การติดตั้ง}{การติดตั้งแบบเฉพาะที่และแบบครอบคลุม}
%-------------------------------------------------------
  ธีมนี้สามารถติดตั้งสำหรับการใช้  \textbf{แบบเฉพาะที่} หรือ \textbf{แบบครอบคลุม} ได้
  \pause
  \begin{block}{การติดตั้งแบบเฉพาะที่}
  \begin{itemize}    
    \item การติดตั้งแบบเฉพาะที่ เป็นวิธีที่ง่ายที่สุดในการติดตั้งธีมนี้ 
    \item คุณต้องเอาไฟล์ต้นกำเนิดทั้ง  4 ไฟล์ไว้ในโฟล์เดอร์เดียวกับไฟล์นำเสนอของคุณ เมื่อคุณดาวน์โหลดธีมนี้ ทั้ง 4 ไฟล์จะอยู่ในโฟล์เดอร์ชื่อ {\tt local}
  \end{itemize}
  \end{block}
  \pause
  \begin{block}{การติดตั้งแบบครอบคลุม}
  \begin{itemize}
     \item ถ้าคุณอยากจะทำให้ธีมนี้ใช้ได้แบบครอบคลุม คุณต้องเอาไฟล์ต้นกำเนิดใส่ไว้ในต้นไม้ไดเร็คทอรี่ของ latex บนเครื่องของคุณ ตำแหน่งที่ตั้งของของ root ของต้นไม้ไดเร็คทอรี่ของ latex บนเครื่องของคุณ จะขึ้นอยู่กับระบบปฎิบัติการและ latex distribution ที่คุณใช้ 
     \item รายละเอียดเกี่ยวกับการติดตั้งภายใต้ระบบปฎิบัติการแบบต่าง ๆ สามารถอ่านได้จากเอกสารของ Beamer
  \end{itemize}
  \end{block}
\end{frame}
     

%-------------------------------------------------------
\subsection{แพ็คเก็จที่ต้องมี}
\begin{frame}{การติดตั้ง}{แพ็คเก็จที่ต้องมี}
%-------------------------------------------------------

  ในการใช้ธีม Feather คุณจะต้องมี  Bemaer class ติดตั้งไว้อยู่แล้ว และต้องมีแพ็คเก็จ 2 แพ็คเก็จ ต่อไปนี้
  \begin{itemize}
    \item TikZ\footnote{TikZ เป็นแพ็คเก็จสำหรับใช้สร้างกราฟฟิกส์ที่สวยงาม ลองดู\chref{http://www.texample.net/tikz/examples/}{ตัวอย่างออนไลน์เหล่านี้} หรือ  \chref{http://tug.ctan.org/tex-archive/graphics/pgf/base/doc/generic/pgf/pgfmanual.pdf}{คู่มือการใช้ pgf}}
    \item calc
  \end{itemize}
  เนื่องจากว่าแพ็คเก็จเหล่านี้เป็นที่ใช้กันแพร่หลาย พวกมันน่าจะมากับ latex distribution ที่คุณใช้ตั้งแต่แรกแล้ว 
\end{frame}

%-------------------------------------------------------
\section{ส่วนต่อประสานผู้ใช้}
\subsection{การเรียกใช้ธีมและตัวเลือกของธีม}
\begin{frame}[shrink]{ส่วนต่อประสานผู้ใช้}{การเรียกใช้ธีมและตัวเลือกของธีม}
%-------------------------------------------------------

  \begin{block}{ธีมการนำเสนอ}
    ธีม Feather สามารถเรียกใช้งานในแบบที่เป็นที่คุ้นเคยอยู่แล้ว  โดยในส่วน  preamble ของไฟล์  {\tt tex} คุณต้องพิมพ์ \\ \vspace{5pt} 
    {\tt \textbackslash usetheme[<options>]\{Feather\}}\\ \vspace{5pt} 
    ธีมการนำเสนอ จะเรียกไฟล์ธีม Feather ชั้นใน ชั้นนอก และสี และส่งต่อ  {\tt <options>} ไปยังไฟล์เหล่านั้น
  \end{block}
  \begin{block}{ธีมชั้นในและชั้นนอก}
    คุณสามารถเรียกใช้แต่ธีมชั้นในหรือชั้นนอกเท่านั้นได้โดยตรง โดยพิมพ์ \\ \vspace{5pt} 
    {\tt \textbackslash useinnertheme\{Feather\}} (และไม่มีตัวเลือก options)\\ \vspace{5pt} 
    {\tt \textbackslash useoutertheme[<options>]\{Feather\}} (ซึ่งมีหนึ่งตัวเลือก option คือ)\\
    \hspace{20pt}{\tt progressstyle=\{fixedCircCnt หรือ } {\tt movingCircCnt\}} \\
    \begin{itemize}
    \item ซึ่งเป็นการตั้งว่าจะให้ความก้าวหน้า progress แสดงอย่างไร
    \item ค่า {\tt movingCircCnt} เป็นค่าปริยาย
    \end{itemize}
  \end{block}
\end{frame}

\begin{frame}[shrink]{ส่วนต่อประสานผู้ใช้}{การเรียกใช้ธีมและตัวเลือกของธีม}

  \begin{block}{ธีมสี}
    นอกจากนี้ คุณสามารถเรียกใช้แต่ธีมสีเท่านั้น โดยเขียนใน preamble ของไฟล์ {\tt tex} ดังนี้ 
    
    \vspace{5pt} 
    
    \begin{itemize}
    \item {\tt \textbackslash usecolortheme\{Feather\}}
    \end{itemize}
    
    \vspace{5pt}
    
    ...หรือจะเปลี่ยนสีขององค์ประกอบต่าง ๆ ในธีม
    
    \vspace{5pt} 
    \begin{itemize}
    \item เปลี่ยนสีของแถบ: \\    
    {\tt \textbackslash setbeamercolor \{Feather\}\{fg=<color>, bg=<color>\}}
    
    \vspace{2pt} 
    
    \item เปลี่ยนสีขององค์ประกอบที่เป็นโครงสร้าง: \\    
    {\tt \textbackslash setbeamercolor\{structure\}\{fg=<color>\}}
    
    \vspace{2pt} 
    
    \item เปลี่ยนสีของตัวอักษรที่เป็นชื่อหรือหัวข้อของเฟรม:\\
    {\tt \textbackslash setbeamercolor\{frametitle\}\{fg=<color>\}}
    
    \vspace{2pt} 
    
    \item เปลี่ยนสีพี้นหลังของตัวอักษรปรกติ:    
    {\tt \textbackslash setbeamercolor\{normal text\}\{fg=<color>, bg=<color>\}}
    \end{itemize}
  \end{block}
\end{frame}


%-------------------------------------------------------
\subsection{รูปขนนก Feather}
\begin{frame}{ส่วนต่อประสานผู้ใช้}{พี้นหลังรูปขนนก Feather}
%-------------------------------------------------------

\begin{block}{พี้นหลังรูปขนนก Feather}
    \begin{itemize}
    \item ในธีม Feather เฟรมหน้าชื่อเรื่อง และ เฟรมหน้าสุดท้าย จะมีรูปขนนก Feather เป็นพื้นหลัง 
    \item คุณสามารถใส่พี้นหลังรูปขนนก Feather ในเฟรมไหนก็ได้โดยการเขียนสิ่งต่อไปนี้ตอนเริ่มต้นของเฟรม
    \end{itemize} 
    
    \vspace{5pt} 
    
  {\tt \{\textbackslash 1bg\\
    \textbackslash begin\{frame\}[<options>]\{Frame Title\}\{Frame Subtitle\}\\
    \ldots\\
    \textbackslash end\{frame\}\}}
\end{block}
\end{frame}

{\1
\begin{frame}[plain,noframenumbering]
  \finalpage{ขอบคุุณที่ใช้บีมเมอร์ธีมขนนก Feather!}
\end{frame}}

\end{document}

\end{document}